\section{第6章の解答}\label{解答:第6章}

\paragraph{6.1}
括弧の中を $t=3x^2-4x+7$ と置くと,\,$\dd t=(6x-4)\dd x$ である.したがって
\[
 \int(6x-4)(3x^2-4x+7)^5\dd x
 =\int t^5\dd t
 =\frac{(3x^2-4x+7)^6}{6}+C.
\]
微分すれば,外側の六乗から $6$ が現れて分母と打ち消し合い,内側の微分 $6x-4$ が掛かって元の式に戻る.多項式なので全実数上で成り立つ.

\paragraph{6.2}
$t=\log x$ と置くと $\dd t=\frac1x\dd x$ だから,
\[
 \int\frac1{x\log x}\dd x
 =\int\frac1t\dd t
 =\log|\log x|+C.
\]
実数の対数のために $x>0$ が必要で,さらに分母の $\log x$ が $0$ となる $x=1$ を除く.よって適用する区間は $0<x<1$ と $1<x$ に分かれ,それぞれの任意定数は独立である.\,$0<x<1$ では $\log x<0$ なので,外側の対数の絶対値を外せない.

\paragraph{6.3}
$t=x^2$,\,$\dd t=2x\dd x$ と置く.\,$x=-2$ は $t=4$,\,$x=1$ は $t=1$ に対応するので,
\[
 \int_{-2}^1 2x\e^{x^2}\dd x
 =\int_4^1\e^t\dd t
 =[\e^t]_4^1=\e-\e^4.
\]
合成形の置換は,\,$\frac{\dd}{\dd x}\e^{x^2}=2x\e^{x^2}$ に基づいており,\,$x^2$ が単調である必要はない.区間を分けても
\[
 \int_4^0\e^t\dd t+\int_0^1\e^t\dd t
 =\int_4^1\e^t\dd t
\]
となる.被積分関数が負になる部分があるため,定積分が負でも矛盾しない.

\paragraph{6.4}
$t=\sqrt{1+x}$ と置けば,\,$x=t^2-1$,\,$\dd x=2t\dd t$ である.端点は $x=0,3$ から $t=1,2$ へ移る.したがって
\[
\begin{aligned}
 \int_0^3\frac1{1+\sqrt{1+x}}\dd x
 &=2\int_1^2\frac{t}{1+t}\dd t\\
 &=2\int_1^2\left(1-\frac1{1+t}\right)\dd t\\
 &=2[t-\log(1+t)]_1^2\\
 &=2-2\log\frac32.
\end{aligned}
\]
$1\le t\le2$ では $1+t>0$ なので,対数の絶対値は不要である.

\paragraph{6.5}
$x=a\sin\theta$ と置き,\,$0\le\theta\le\frac\pi6$ を選ぶ.この区間で $\cos\theta>0$ だから,\,$\sqrt{a^2-x^2}=a\cos\theta$ と書ける.また $\dd x=a\cos\theta\dd\theta$ なので,
\[
\begin{aligned}
 \int_0^{\frac a2}\sqrt{a^2-x^2}\dd x
 &=a^2\int_0^{\frac\pi6}\cos^2\theta\dd\theta\\
 &=a^2\left[\frac\theta2+\frac14\sin(2\theta)\right]_0^{\frac\pi6}\\
 &=a^2\left(\frac\pi{12}+\frac{\sqrt3}{8}\right).
\end{aligned}
\]
最後は $\sin\frac\pi3=\frac{\sqrt3}{2}$ を用いた.

\paragraph{6.6}
$0\le\theta\le\frac\pi3$ で $x=\tan\theta$ と置く.端点は $0,\sqrt3$ に対応し,この区間では $\cos\theta\ne0$ である.よって
\[
 \int_0^{\sqrt3}\frac1{1+x^2}\dd x
 =\int_0^{\frac\pi3}\frac1{1+\tan^2\theta}
   \frac1{\cos^2\theta}\dd\theta
 =\int_0^{\frac\pi3}1\dd\theta=\frac\pi3.
\]
分母と微分に現れる因子が等しいことを使っている.

\paragraph{6.7}
変数を換えるときは,端点だけでなく被積分関数と微分の記号も書き換える必要がある.\,$t=x^2$ では $\dd t=2x\dd x$ であり,\,$\dd x$ をそのまま $\dd t$ にすることはできない.また $-1\le x\le1$ 全体で $x=\sqrt{t}$ と戻すこともできない.負の $x$ は $x=-\sqrt{t}$ という別の枝に属する.したがって,\,$1$ の積分を同じ端点の $t$ の積分として表す根拠がない.実際には $\int_{-1}^1 1\dd x=2$ である.

一方,第二の積分には $2x\dd x$ があるから,合成形の置換をそのまま使える.
\[
 \int_{-1}^1\frac{2x}{1+x^2}\dd x
 =\int_1^1\frac1{1+t}\dd t=0.
\]
原始関数 $\log(1+x^2)$ の上下端の値がともに $\log2$ になることからも確かめられる.同じ端点へ移ることだけでなく,置換公式の形を満たしていることが結論を支えている.
